\oppgave{8}
Vi tolker bipolar NRZ i oppgaven som antipodal, binær signalering med to motsatte spenningsnivåer, altså $-A$ og $+A$. Vi antar like sannsynlige bit, additiv hvit Gauss-støy og optimal deteksjon uten intersymbolinterferens. Bitfeilsannsynligheten er da:
\[
    \mathrm{BER}=P_b=Q\!\left(\sqrt{\frac{2E_b}{N_0}}\right)
\]
$Q$-funksjonen gir sannsynligheten for at en standard normalfordelt variabel er større enn argumentet.

\subsection*{a)}
Kravet til bitfeilsannsynlighet er:
\[
    Q\!\left(\sqrt{\frac{2E_b}{N_0}}\right)\leq10^{-6}
\]
Siden $Q(x)$ er avtakende, finner vi minste nødvendige bitenergiforhold ved å sette BER lik grenseverdien:
\[
    \sqrt{\frac{2E_b}{N_0}}=Q^{-1}(10^{-6})
\]
Verdien av $Q^{-1}(10^{-6})$ må finnes fra en normalfordelingstabell eller en numerisk funksjon. Den inverse funksjonen finner argumentet som gir den ønskede halesannsynligheten, og er ikke det samme som $1/Q$.

En beregning uten ekstra Matlab/Octave-toolbox bruker sammenhengen med den komplementære feilfunksjonen og dens inverse:
\begin{align*}
    Q(x)&=\frac{1}{2}\operatorname{erfc}\!\left(\frac{x}{\sqrt{2}}\right)\\
    Q^{-1}(p)&=\sqrt{2}\operatorname{erfc}^{-1}(2p)
\end{align*}
For det oppgitte BER-kravet blir beregningen:
\[
    Q^{-1}(10^{-6})=\sqrt{2}\operatorname{erfc}^{-1}(2\cdot10^{-6})
                  \approx4.753
\]
Den samme tallverdien beregnes i Matlab/Octave med:
\begin{verbatim}
x = sqrt(2)*erfcinv(2*1e-6);
\end{verbatim}
Vi kvadrerer og deler på 2 for å finne den lineære verdien:
\[
    \left(\frac{E_b}{N_0}\right)_{\min}
       =\frac{(4.753)^2}{2}\approx11.298
\]
Dette omregnes til dB:
\[
    \left(\frac{E_b}{N_0}\right)_{\min,\mathrm{dB}}
       =10\log_{10}(11.298)\approx10.530\enhet{dB}
\]
Minste verdi som oppfyller BER-kravet er derfor:
\[
    \boxed{\left(\frac{E_b}{N_0}\right)_{\min}\approx11.298,
    \quad\left(\frac{E_b}{N_0}\right)_{\min,\mathrm{dB}}\approx10.530\enhet{dB}}
\]

\newpage
\subsection*{b)}
Det oppgitte bitenergiforholdet er:
\[
    \left(\frac{E_b}{N_0}\right)_{\mathrm{dB}}=11\enhet{dB}
\]
Først omregnes det til lineær verdi:
\[
    \frac{E_b}{N_0}=10^{11/10}\approx12.589
\]
Dette settes inn i BER-uttrykket for antipodal NRZ:
\[
    P_b=Q\!\left(\sqrt{2\cdot12.589}\right)
\]
Argumentet til $Q$-funksjonen blir:
\[
    x=\sqrt{25.178}\approx5.018
\]
Ved å bruke den komplementære feilfunksjonen får vi:
\[
    P_b=\frac{1}{2}\operatorname{erfc}\!\left(    \frac{5.018}{\sqrt{2}}\right)
       \approx2.613\cdot10^{-7}
\]
Dette kan beregnes direkte i Matlab/Octave med:
\begin{verbatim}
Y = 10^(11/10);
BER = 0.5*erfc(sqrt(Y));
\end{verbatim}
Den forventede bitfeilsannsynligheten er dermed:
\[
    \boxed{\mathrm{BER}\approx2.613\cdot10^{-7}}
\]
Dette er lavere enn $10^{-6}$, i samsvar med at $11\enhet{dB}$ er høyere enn den nødvendige minsteverdien fra a).
